
\part{Сесія 03: джерельні модулі, вимірювання, керування і фільтри}

\chapter{Архітектура: зведений довідник і окремі статті}

Цей документ є \emph{зведеним інженерно-математичним ядром}. Це не випадкова компіляція: він зводить кілька джерельних ліній в одну робочу систему позначень, бо користувачеві потрібна не лише бібліографія, а швидко доступна математична інфраструктура. Окремі статті й книги залишаються корисними, але вони мають різні стилі, різні домовленості щодо кватерніонів, різні позначення для збурень, різні одиниці виміру шуму і різний рівень практичної деталізації. Зведений довідник знімає цю несумісність.

\section{Два паралельні продукти}

Надалі треба вести дві гілки.

\begin{enumerate}
\item \textbf{Довідникова гілка.} Це поточний файл: український TeX-довідник з єдиними позначеннями, короткими поясненнями, контрольними формулами і мінімальними алгоритмічними шаблонами. Він придатний для швидкого навчання, рецензування, пошуку помилок і машинного добудовування.
\item \textbf{Джерельна гілка.} Для самодостатніх статей з arXiv або відкритих репозиторіїв робиться окремий каталог, де структура оригіналу збережена: назва, розділи, файли, бібліографія, ліцензійні й авторські нотатки. Там переклад має бути ближчим до оригіналу.
\end{enumerate}

Практична різниця така. Якщо треба швидко зрозуміти, як перевірити коректність ESKF, користувач відкриває довідникову гілку. Якщо треба цитувати або звіряти конкретну статтю, користувач відкриває джерельну гілку.

\section{Що вважається ``першим ядром''}

Перший пріоритет має не весь корпус, а мінімальний набір, без якого решта стає крихкою:

\begin{enumerate}
\item фільтрація стану: Байєсівський крок, KF, EKF, ESKF, коваріації, NIS/NEES;
\item обертання і групи Лі: $\SO(3)$, $S^3$, $\SE(3)$, $\Exp$, $\Log$, $\Ad$, ліві/праві збурення;
\item IMU-моделі: номінальний стан, похибковий стан, зміщення, шум, неперервно-дискретна коваріація;
\item сигнали: згортка, спектр, дискретизація, FIR/IIR, вікна, I/Q-представлення;
\item керування: простір станів, спостережуваність, керованість, PID, LQR, насичення, стійкість;
\item чисельна математика: лінійні розв'язувачі, найменші квадрати, регуляризація, скінченні різниці, перевірки збіжності.
\end{enumerate}

\section{Правило для окремих статей}

Якщо джерело є одиничною статтею з чистим TeX, то воно отримує власний підпакет. Мінімальний підпакет має містити:

\begin{itemize}
\item \texttt{main.tex} з компільованим українським перекладом або перекладеним стартом;
\item \texttt{SOURCE.md} з артикулом, авторами, ідентифікатором, станом ліцензії й межами перекладу;
\item \texttt{translation\_status.csv} з переліком розділів: \texttt{not-started}, \texttt{draft-ua}, \texttt{formula-checked}, \texttt{reviewed};
\item збережені ключові макроси або мапа макросів, якщо оригінальний TeX має власні команди;
\item окремий PDF для перегляду.
\end{itemize}

\section{Правило для великих книг}

Великі книги не треба перекладати лінійно з першої сторінки. Їх треба різати на практичні модулі. Наприклад, для книги з робототехніки спочатку беруться координатні перетворення, кінематика, локалізація, базове керування, сенсори і планування. Історичні, педагогічні або надто повільні розділи йдуть пізніше. Для чисельних PDE першими йдуть сітки, стійкість, CFL, хвильове рівняння, перевірка збіжності й кодові тести.

\section{Машинозчитуваність як технічна вимога}

Кожен модуль повинен компілюватися або принаймні проходити TeX-парсинг. Формули мають бути структурованими, а не збереженими як зображення. Позначення повинні бути повторно використовними. Неперевірені місця позначаються явно:

\[
\texttt{[[CHECK\_FORMULA]]},\qquad
\texttt{[[CHECK\_TERM]]},\qquad
\texttt{[[SOURCE\_ALIGNMENT\_NEEDED]]}.
\]

Такі мітки не є провалом. Вони роблять стан роботи явним і придатним для локальних агентів.
